Fix \(0<e<1\) and set \(c=1/e\).  If \(P\ll Q\) and \(r=dP/dQ\), then \(\int r\,dQ=1\).  Therefore replacing \(f\) by \(f-b(t-1)\) changes neither \(D_f(P\|Q)\) nor \(B_f(e)\).  Convexity gives the chord bound
\[
 f(t)\le \left(1-\frac{t}{c}\right)f(0)+\frac{t}{c}f(c),\qquad 0\le t\le c. \tag{1}
\]
If \(Q=eP+(1-e)R\), then \(Q\ge eP\), so \(P\ll Q\) and \(r\le c\), \(Q\)-almost surely.  Integrating (1) and using \(\int r\,dQ=1\) yields
\[
 D_f(P\|Q)\le (1-e)f(0)+ef(c)=B_f(e). \tag{2}
\]
Thus \(\mathfrak L^{\to}_e(P)\subseteq\mathfrak B^{\to}_{f,e}(P)\) for every admissible generator.

Suppose \(f\in\mathfrak H_c\).  Choose \(b\) such that \(g(t)=f(t)-b(t-1)\) vanishes on \([0,c]\) and is strictly positive on \((c,\infty)\).  Affine invariance gives \(B_g(e)=B_f(e)=0\).  If \(Q\in\mathfrak B^{\to}_{f,e}(P)\), then
\[
 0\le\int g(r)\,dQ=D_f(P\|Q)\le B_f(e)=0.
\]
Consequently \(r\le c\), \(Q\)-almost surely.  The measure \((Q-eP)/(1-e)\) has nonnegative \(Q\)-density \((1-er)/(1-e)\) and total mass one, so it is a probability law.  Hence \(Q\in\mathfrak L^{\to}_e(P)\), and equality follows from (2).

For necessity, assume equality for every full-support binary \(P\).  Let \(P=\operatorname{Bernoulli}(p)\), \(0<p<1\), and \(Q_q=\operatorname{Bernoulli}(q)\).  The boundary law \(Q_{ep}\) is in \(\mathfrak L^{\to}_e(P)\), whereas \(0<q<ep\) gives \(p/q>c\) and hence \(Q_q\notin\mathfrak L^{\to}_e(P)\).  Exactness implies
\[
 D_f(P\|Q_q)>B_f(e)\quad(0<q<ep),
 \qquad D_f(P\|Q_{ep})\le B_f(e).
\]
Because \(f\) is finite and continuous and \(q\in(0,1)\), the binary divergence is continuous in \(q\).  Letting \(q\uparrow ep\) proves equality at the boundary:
\[
 ep f(c)+(1-ep)f(t_p)=B_f(e),
 \qquad t_p=\frac{1-p}{1-ep}. \tag{3}
\]
Letting \(p\downarrow0\), using \(t_p\to1\) and \(f(1)=0\), gives \(B_f(e)=0\).

Define \(g(t)=f(t)+f(0)(t-1)\).  The radius identity is \(f(c)+(c-1)f(0)=0\), and therefore
\[
 g(0)=g(1)=g(c)=0. \tag{4}
\]
Affine invariance and (3) give \(g(t_p)=0\) for every \(p\in(0,1)\).  Since \(p\mapsto t_p\) ranges from \(1\) down to \(0\), continuity gives \(g=0\) on \([0,1]\).

Now take \(y\in(1,c)\), \(x\in(0,1)\), and \(\lambda=(1-x)/(y-x)\).  On a binary alphabet let \(Q\) assign weights \(\lambda,1-\lambda\), and define \(P\) by likelihood ratios \(y,x\).  Since
\[
 \lambda y+(1-\lambda)x=1,
\]
these are full-support probability laws, and their likelihood ratio is bounded by \(c\).  The cap characterization places \(Q\) in \(\mathfrak L^{\to}_e(P)\).  Exactness, (4), and convexity at the mean give
\[
 0=g(1)\le D_g(P\|Q)=\lambda g(y)+(1-\lambda)g(x)\le B_g(e)=0.
\]
As \(g(x)=0\), one gets \(g(y)=0\).  Hence \(g=0\) on \([0,c]\).

For \(t>c\), convexity applied with \(c\) between \(1\) and \(t\) gives \(g(t)\ge0\).  If \(g(t)=0\), repeat the two-ratio construction with ratios \(t\) and any \(x\in(0,1)\).  It gives \(D_g(P\|Q)=0=B_g(e)\), so \(Q\) is in the divergence ball, but its likelihood ratio exceeds \(c\), so it is not in the mixture class.  This contradicts binary exactness.  Thus \(g(t)>0\) for every \(t>c\), which is precisely \(f\in\mathfrak H_c\).  Universal equality implies binary equality trivially, and the preceding arguments prove the reverse implications.  The one-sided bow-mixture theorem supplies the causal interpretation of the mixture class.